\section*{Chapitre \ref{ch:ml:deep-hedging-and-machine-learning-in-pricing} --- Couverture profonde et apprentissage automatique dans la tarification}

\begin{solution}{exo:ml:deep-hedging-and-machine-learning-in-pricing:1}
$\rho(X + c) = \lambda^{-1}\log\E[e^{-\lambda X}e^{-\lambda c}] = \rho(X) - c$. Pour une constante, $\rho(c) = -c$~: détenir la trésorerie $c$
exige $-c$ de trésorerie supplémentaire pour être acceptable, c'est-à-dire qu'elle libère $c$.
\end{solution}

\begin{solution}{exo:ml:deep-hedging-and-machine-learning-in-pricing:2}
$(1.5\times0.001\times100\times0.0998^2)^{1/3} = (0.001494)^{1/3} = 0.114$~: une bande $\pm0.114$ en delta, large de 0.229, contre la
bande apprise de 0.055, au même coût.
\end{solution}

\begin{solution}{exo:ml:deep-hedging-and-machine-learning-in-pricing:3}
La couverture quotidienne avec un seul instrument ne peut pas éliminer le risque de rééquilibrage discret ni celui de la volatilité stochastique, et le vendeur
facture le risque résiduel avec $\lambda = 1$~: 2.454 contre 2.254, soit une prime de risque de 0.20.
\end{solution}

\begin{solution}{exo:ml:deep-hedging-and-machine-learning-in-pricing:4}
Sous Heston avec une corrélation négative, une baisse du spot s'accompagne d'une hausse de la volatilité, ce qui augmente la valeur du call~; le
ratio de couverture à variance minimale est le delta corrigé d'un terme tenant compte de ce co-mouvement (couverture plus faible pour un call lorsque $\rho <
0$). Le delta de Black--Scholes, à volatilité implicite fixée, l'ignore~; le réseau, entraîné sur des trajectoires de Heston, l'apprend.
\end{solution}

\begin{solution}{exo:ml:deep-hedging-and-machine-learning-in-pricing:5}
Leurs effets sur les prix sont presque interchangeables sur trois échéances~: un retour à la moyenne plus rapide vers un niveau de long terme et
un retour plus lent vers un autre niveau peuvent produire des structures par terme semblables. Des échéances plus longues, davantage d'échéances et
des instruments sensibles à la variance (swaps de variance) les départagent.
\end{solution}

\begin{solution}{exo:ml:deep-hedging-and-machine-learning-in-pricing:6}
Les scénarios de stress se situent hors de la région échantillonnée par l'\omterm{def:ml:why-financial-machine-learning-is-different:sets}{ensemble de test}, là où le \omterm{def:ml:deep-hedging-and-machine-learning-in-pricing:surrogate}{modèle de substitution} extrapole~; son erreur de test n'en dit rien
là-bas. Valider sur les scénarios eux-mêmes, ou entraîner sur un pavé qui les contient, et vérifier les grecques comme les prix.
\end{solution}

\begin{solution}{exo:ml:deep-hedging-and-machine-learning-in-pricing:7}
La dérivée trajectorielle de $\mathbf1\{S_T > K\}$ par rapport à $S_0$ est presque sûrement nulle~: l'étiquette indique que le delta est nul
alors que le vrai delta est une densité au strike. Les étiquettes trajectorielles exigent un payoff régulier~; le remède usuel consiste à lisser la
discontinuité (remplacer le digital par un call spread serré) avant de dériver, ou utiliser des étiquettes fondées sur le rapport de vraisemblance.
\end{solution}

\begin{solution}{exo:ml:deep-hedging-and-machine-learning-in-pricing:8}
Sur un marché complet, le produit est répliqué exactement~: une stratégie $H$ a pour gains $G_H = Z - p_{BS}$ trajectoire par trajectoire. Toute stratégie
pour le vendeur peut s'écrire $H + \delta'$, donc $-Z + G_{H + \delta'} = -p_{BS} + G_{\delta'}$ et, par invariance à la trésorerie,
$\rho(-Z + G_{H+\delta'}) = p_{BS} + \rho(G_{\delta'})$. En prenant l'infimum sur $\delta'$ et en retranchant le terme de référence sans produit
$\inf_{\delta'}\rho(G_{\delta'})$ laisse $p = p_{BS}$, quel que soit $\lambda$.
\end{solution}

\begin{solution}{pb:ml:deep-hedging-and-machine-learning-in-pricing:1}
\textbf{Partie I.}
\begin{enumerate}
\item Un call à la monnaie d'un mois sur 100 sous Heston, couvert quotidiennement avec le sous-jacent~; le réseau voit le temps jusqu'à
l'échéance, la log-moneyness et, par la bande, sa position courante.
\item 0.123 en moyenne, soit plus d'un point de volatilité de vega (0.115).
\item Ne pas trader est alors le choix par défaut, la bande est lisible, et l'entraînement est plus aisé que d'apprendre à reproduire
exactement la position précédente.
\item Le risque entropique du P\&L avant la prime~; $\lambda = 1$ mesure une aversion absolue au risque par unité de prix, donc une
perte de 1 pèse $e$ fois autant qu'une perte de 0.
\end{enumerate}
\textbf{Partie II.}
\begin{enumerate}[resume]
\item À 0, 5, 10, 20 points de base~: deep 2.454, 2.558, 2.658, 2.840~; delta de Black--Scholes 2.475, 2.608, 2.743, 3.018~;
Whalley--Wilmott 2.475, 2.606, 2.712, 2.910.
\item À la monnaie, à mi-chemin de l'échéance, la bande apprise a une largeur de 0.031 à 5 points de base et de 0.104 à 20, contre 0.181 et
0.288~; elle est plus étroite à la monnaie et au-dessus, et plus large en dessous de la monnaie.
\item L'expected shortfall est plus faible (3.763 contre 3.817) et le risque entropique plus élevé (2.671 contre 2.658)~: chaque objectif l'emporte selon
sa propre mesure.
\item Sa bande large trade le moins (coût moyen 2.424), mais la position s'écarte davantage de la couverture, et l'écart-type
du P\&L (0.703) et les queues sont les plus fortes parmi les couvertures.
\end{enumerate}
\textbf{Partie III.}
\begin{enumerate}[resume]
\item Précision par échantillon~: avec 1\,024 échantillons, erreur sur le prix de 0.218 contre 0.436 et erreur sur le delta de 0.021 contre 0.059~; à 4\,096,
prix égaux et deltas toujours meilleurs.
\item Erreurs sur les paramètres de 1.7~\% (variance initiale) à 18.7~\% (vitesse de retour à la moyenne) de chaque plage, 13.2 points de base sur le spot dans
les surfaces repricées, 0.46 points de volatilité implicite sur cinq surfaces.
\item La réponse du réseau sert de point de départ à l'optimiseur (et de contrôle de son résultat).
\item Hors du pavé d'entraînement, sur les contraintes d'arbitrage, et sur les données de marché qu'aucun modèle n'ajuste.
\end{enumerate}
\textbf{Partie IV.}
\begin{enumerate}[resume]
\item \emph{Couverture avec frictions.} Le \omterm{def:ml:deep-hedging-and-machine-learning-in-pricing:risk}{prix d'indifférence} passe de 2.454 sans coûts à 2.558, 2.658 et 2.840 à 5, 10
et 20 points de base~; la bande apprise à la monnaie a une largeur de 0.031 et 0.104 à 5 et 20 points de base, contre celle de Whalley--Wilmott
0.181 et 0.288.
\item Comme borne sur le prix auquel le desk est rémunéré pour le risque qu'il conserve, et non comme cotation~: le prix de marché et
les autres positions du portefeuille comptent.
\item Trajectoires hors échantillon (out-of-sample), autres modèles et autres paramètres, trajectoires de stress, et comparaison avec la bande asymptotique et le delta
sur l'historique du portefeuille.
\item Le réseau choisirait aussi la quantité du second instrument à détenir pour couvrir le risque de volatilité~; le \omterm{def:ml:deep-hedging-and-machine-learning-in-pricing:risk}{prix d'indifférence}
baisserait vers la valeur de Heston.
\item Réentraîner le réseau ou le confronter au pricer de référence après chaque recalibration, sur une grille couvrant les nouveaux paramètres.
\item La couverture est optimale pour une dynamique de volatilité~; une autre (volatilité de la variance plus élevée, saut) peut la rendre
pire que le delta.
\item Elle calcule les dérivées propres du réseau en \omterm{def:ml:deep-hedging-and-machine-learning-in-pricing:surrogate}{apprentissage automatique différentiel}, les gradients de la \omterm{def:ml:deep-hedging-and-machine-learning-in-pricing:dh}{couverture profonde}, et les
grecques exactes du pricer de référence.
\item Quand le trading vaut son coût.
\end{enumerate}
\end{solution}

\begin{solution}{iq:ml:deep-hedging-and-machine-learning-in-pricing:1}
Chaque rééquilibrage a un coût~; de petits écarts par rapport au delta ajoutent peu de risque, alors que les éliminer en trading coûte à chaque fois une fourchette. La
\omterm{def:ml:reinforcement-learning-foundations:mdp}{politique} optimale laisse la position inchangée tant qu'elle est assez proche du delta et la ramène au bord de la bande sinon~; la
bande s'élargit avec les coûts et se resserre avec le gamma et l'aversion au risque.
\iqlookfor{le compromis coût-risque et la manière dont la bande dépend du coût, du gamma et de l'aversion au risque.}
\end{solution}

\begin{solution}{iq:ml:deep-hedging-and-machine-learning-in-pricing:2}
Une \omterm{def:ml:reinforcement-learning-foundations:mdp}{politique} de \omterm{def:ml:neural-networks-for-noisy-tabular-data:mlp}{réseau de neurones} entraînée sur des trajectoires simulées~: les entrées sont l'état observable (temps, prix, positions, autres variables), les sorties sont
les positions dans les instruments de couverture, la perte est une \omterm{def:ml:deep-hedging-and-machine-learning-in-pricing:risk}{mesure de risque convexe} du P\&L terminal couvert, après coûts.
\iqlookfor{état, action, perte et simulation comme données d'entraînement.}
\end{solution}

\begin{solution}{iq:ml:deep-hedging-and-machine-learning-in-pricing:3}
Entraîner sur des étiquettes et leurs dérivées par rapport aux entrées, en faisant coïncider les dérivées du réseau par
différentiation automatique~; les étiquettes de dérivées trajectorielles issues de Monte-Carlo sont presque gratuites et bien moins bruitées relativement à leur signal, de sorte que
moins d'échantillons suffisent à déterminer la fonction et sa pente.
\iqlookfor{la perte et la raison pour laquelle les étiquettes dérivées réduisent le besoin en échantillons.}
\end{solution}

\begin{solution}{iq:ml:deep-hedging-and-machine-learning-in-pricing:4}
La prime qui laisse le vendeur indifférent, au sens d'une fonction d'utilité ou d'une mesure de risque, entre vendre avec la meilleure couverture et ne pas
vendre. Avec une mesure de risque invariante à la trésorerie, c'est le risque de la meilleure position couverte sans la prime.
\iqlookfor{la définition et son lien avec l'invariance à la trésorerie.}
\end{solution}

\begin{solution}{iq:ml:deep-hedging-and-machine-learning-in-pricing:5}
Il est approximatif, aveugle hors de son pavé d'entraînement, entraîné sur des données de modèle qui s'ajustent parfaitement alors que les données de marché ne le font pas, et
ne fournit aucun diagnostic d'ajustement~; il sert à démarrer l'optimiseur et à vérifier son résultat.
\iqlookfor{approximation, extrapolation et décalage entre données de modèle et données de marché.}
\end{solution}

\begin{solution}{iq:ml:deep-hedging-and-machine-learning-in-pricing:6}
Le \omterm{def:ml:deep-hedging-and-machine-learning-in-pricing:surrogate}{modèle de substitution} extrapole ou est sous-entraîné là où les prix sont minuscules, et son ajustement est jugé à l'aune de l'erreur absolue~; rien n'impose
la convexité. Entraîner sur les prix en logarithme ou les volatilités implicites, utiliser des étiquettes différentielles incluant les dérivées secondes, ajouter des contraintes de
convexité ou des pénalités, et échantillonner plus densément les ailes.
\iqlookfor{la cause et au moins deux remèdes.}
\end{solution}
